\section*{Заключение}
\addcontentsline{toc}{section}{Заключение}

В результате проведённого исследования подтверждена гипотеза о том, что
применение элементов геймификации в системах управления ирригацией позволяет
повысить вовлечённость операторов и улучшить контроль за режимом орошения
тропических культур. Разработанная имитационная модель продемонстрировала,
что задача автоматизации полива на плантациях личи требует учёта
микроклиматических особенностей и специфики тропического водного баланса.

На основе полученных при разработке данных можно сформулировать следующие
\textbf{выводы}:

\begin{enumerate}
    \item Использование WebSocket-протокола и архитектуры на базе Go
    (фоновый движок симуляции) и React (клиентский интерфейс)
    обеспечивает необходимую гибкость для потоковой передачи
    результатов имитационной модели в реальном времени без зависимости
    от проприетарных платных сервисов и физических датчиков.
    \item Проектирование модели данных с поддержкой временных рядов
    телеметрии является критическим условием для масштабируемости
    системы, так как позволяет эффективно хранить и анализировать
    историю показаний, генерируемых движком симуляции.
    \item Внедрение игровых механик (виртуальная лейка, индексы здоровья
    растений, система уведомлений о засухе и переувлажнении)
    демонстрирует положительный эффект на регулярность и своевременность
    проведения агротехнических операций в рамках имитационной среды.
\end{enumerate}

\textbf{Рекомендации и предложения.} Для практического применения разработанного
решения в деятельности малых фермерских хозяйств предлагается использовать
модель облачного развёртывания с адаптивным мобильным интерфейсом. Это обеспечит
бесперебойный доступ к системе мониторинга в полевых условиях при минимальных
затратах на ИТ-инфраструктуру.

\textbf{Перспективы углубления темы.} В дальнейшем исследовании планируется
интеграция с реальными IoT-датчиками влажности и температуры (через протокол
MQTT), внедрение модуля предиктивной аналитики для прогнозирования потребности
в поливе на основе метеорологических данных, а также расширение геймификации
за счёт системы достижений и рейтинга операторов.

\textbf{Доступ к исходному коду.} Исходный код разработанного веб-приложения ---
серверная часть на Go, клиентское приложение на React, конфигурация
контейнеризации и SQL-миграции --- опубликован в открытом репозитории GitHub
и доступен по адресу: \url{https://github.com/ZertGraf/CP}.

\textbf{Демонстрационный доступ.} Развёрнутый экземпляр разработанного
веб-приложения доступен для ознакомления и тестирования по адресу:
\url{https://lichi.89-169-170-211.sslip.io}.

При разработке пояснительной записки к курсовому проекту строго соблюдались
требования, определяющие структуру и правила оформления
научно-исследовательских работ.